\documentclass{article}
\usepackage{graphicx} % Required for inserting images

\title{EJERCICIO}
\author{Rafael Josue Gallardo Blanco}
\date{March 2026}

\begin{document}

\maketitle

\section{PROBLEM INDUSTRIAL}
\large
Una panaderia produce Panes y Tortas cada uno con el siguiente margen de ganancia:

\begin{itemize}
    \item Pan : \$2 
    \item Torta : \$6
\end{itemize}

\large El horno tiene una disponibilidad de 18 horas y los tiempos de cada uno son:

\begin{itemize}
    \item Pan : 1 Hora 
    \item Torta : 3 Horas
\end{itemize}

\large la demanda minima de la panaderia es de 2 tortas al dia, ¿cuanto produce para maximizar la ganancia?

\begin{center}
    \item \textbf{Variable XA:} Cantidad de produccion de panes 
    \item \textbf{Variable XB:} Cantidad de produccion de tortas
    \item \textbf{Funcion objetivo:} Maximizar Z: 2XA + 6XB
    \item \large \textbf{Restricciones:} \item 1XA + 3XB = 18 y XB $ \geq 2$, XA$ \geq 0$
    
\end{center}
\vspace{1cm}
\section{PROBLEM SISTEMAS}
\large
 Un balanceador distribuye entre 3 microservicios: \textbf{M1}, \textbf{M2} Y \textbf{M3}
\ cada uno de los microservicios tiene el siguiente costo

\begin{itemize}
    \item M1 : \$0.02
    \item M2 : \$0.015
    \item M3 : \$0.03
\end{itemize}

\large acompañado de las siguientes capacidades maximas:

\begin{itemize}
    \item M1  $ \leq 200$
    \item M2  $ \leq 250$
    \item M3  $ \leq 150$
\end{itemize}

\large El total requerido es de 500 repeticiones por minuto ¿como minimizar el costo total?

\begin{center}
    \item \textbf{Variable XA:} Cantidad de repeticiones en el microservidor M1
    \item \textbf{Variable XB:} Cantidad de repeticiones en el microservidor M2
    \item \textbf{Variable XC:} Cantidad de repeticiones en el microservidor M3
    \item \textbf{Funcion objetivo:} Minimizar Z: 0.02XA + 0.015XB + 0.03XC
    \item \large \textbf{Restricciones:} 200XA + 250XB + 150XC = 500 y XA  $ \geq 0$, XB  $ \geq 0$, XC  $ \geq 0$
\end{center}


\vspace{5cm}
\section{PROBLEMA IDEAL}
\large
 una empresa de lacteos tiene 3 productos \textbf{yogurt}, \textbf{queso} Y \textbf{Leche}
\ cada uno de los productos tienen los siguientes margenes de ganancia

\begin{itemize}
    \item yogurt : \$2000
    \item queso : \$1800
    \item leche : \$2100
\end{itemize}

\large cada producto ocupa las siguientes unidades de espacio por caja: 

\begin{itemize}
    \item yogurt : 2.2
    \item queso : 1.3
    \item leche : 2.5
\end{itemize}

\large la capacidad maximas de las cajas son de 30 unidades ¿como debe ser la distribucuin en la caja para maximizar la ganancia?

\begin{center}
    \item \textbf{Variable XA:} cantidad de yogurt en una caja
    \item \textbf{Variable XB:} cantidad de queso en una caja
    \item \textbf{Variable XC:} cantidad de leche en una caja
    \item \textbf{Funcion objetivo:} Miximizar Z: 2000XA + 1800XB + 2100XC
    \item \large \textbf{Restricciones:} 2.2AX + 1.3XB + 2.5 XC $ \leq 30$ y XA  $ \geq 0$, XB  $ \geq 0$, XC  $ \geq 0$
\end{center}

\end{document}

